\DocumentMetadata{
  pdfversion=2.0,
  pdfstandard=ua-2,
  lang=en-US,
  tagging=on,
  tagging-setup={math/setup=mathml-SE,tabsorder=structure}
}
\documentclass[11pt,letterpaper]{article}
\usepackage[margin=0.68in,includefoot]{geometry}
\usepackage{newcomputermodern}
\usepackage{amsmath,amscd,mathtools}
\usepackage{tikz,adjustbox}
\providecommand{\square}{\mdlgwhtsquare}
\usepackage[hidelinks]{hyperref}
\hypersetup{
  pdftitle={Topology First-Year Exam, January 2025, Part 2},
  pdfauthor={University of Florida Department of Mathematics},
  pdfsubject={Graduate mathematics examination},
  pdfdisplaydoctitle=true,
  bookmarksopen=true
}
\setcounter{secnumdepth}{0}
\setlength{\parindent}{0pt}
\setlength{\parskip}{0.28em}
\setlength{\emergencystretch}{3em}
\setlength{\leftmargini}{2.15em}
\setlength{\leftmarginii}{2.65em}
\setlength{\leftmarginiii}{3em}
\renewcommand{\labelenumi}{\textbf{\theenumi.}}
\pagestyle{empty}
\raggedbottom
\allowdisplaybreaks
\newenvironment{examproblems}[1]{%
  \begin{enumerate}%
  \setcounter{enumi}{#1}%
  \setlength{\topsep}{0.35em}%
  \setlength{\partopsep}{0pt}%
  \setlength{\itemsep}{0.48em}%
  \setlength{\parsep}{0.18em}%
}{\end{enumerate}}
\newenvironment{examsubparts}{%
  \begin{enumerate}%
  \setlength{\topsep}{0.18em}%
  \setlength{\partopsep}{0pt}%
  \setlength{\itemsep}{0.16em}%
  \setlength{\parsep}{0.1em}%
}{\end{enumerate}}
\newenvironment{examinstructions}{%
  \par\begingroup\itshape
}{%
  \par\endgroup\vspace{0.18em}
}
\makeatletter
\renewcommand\section{\@startsection{section}{1}{0pt}%
  {-1.25ex plus -.25ex minus -.1ex}%
  {0.55ex plus .1ex}%
  {\normalfont\large\bfseries}}
\renewcommand{\@maketitle}{%
  \newpage\null\vskip -1em%
  {\footnotesize\bfseries
    \href{https://www.ufl.edu/}{University of Florida}, \href{https://math.ufl.edu/}{Department of Mathematics}\par}%
  \vskip -0.5em%
  \tagstructbegin{tag=Title}%
    {\LARGE\bfseries\href{https://gma.math.ufl.edu/exams/topology-first-year/}{Topology First-Year Exam}, January 2025, Part 2\par}%
  \tagstructend%
  \vskip 0.25em%
  \tagmcbegin{artifact}%
    \hrule height 1pt%
  \tagmcend%
  \vskip 1em%
}
\makeatother
\title{Topology First-Year Exam, January 2025, Part 2}
\author{}
\date{}
\begin{document}
\maketitle
\thispagestyle{empty}
\begin{examinstructions}
For the first five problems, show all your work and support all statements.
\end{examinstructions}
\begin{examproblems}{0}
\item
Suppose that the continuous maps \(f,g:X\to Y\) are homotopic and that the continuous maps \(h,k:Y\to Z\) are homotopic. Show that \(h\circ f\) is homotopic to \(k\circ g\).
\item
Prove that \(\pi_1(S^1,1)\cong\mathbb{Z}\).
\item
State and prove the Brouwer fixed point theorem for the unit disk, \(D^2\).
\item
Prove that for \(n\geq 2\), the fundamental group of the~\(n\)-sphere, \(S^n\), is~\(0\).
\item
Show that the torus \(T=S^1\times S^1\) is not homeomorphic to the surface of genus two, \(M_2=T\mathbin{\#}T\).
\end{examproblems}
\begin{examinstructions}
Answer the following with complete definitions or statements or short proofs.
\end{examinstructions}
\begin{examproblems}{5}
\item
Suppose that \(f,g:X\to\mathbb{R}^n\) are continuous. Show that~\(f\) and~\(g\) are homotopic.
\item
State the Borsuk–Ulam Theorem for \(S^2\).
\item
State the Seifert–van Kampen Theorem.
\item
State the Tietze Extension Theorem.
\item
Show that \(\mathbb{R}^2\) is not homeomorphic to \(\mathbb{R}^3\).
\item
State the Tychonoff theorem.
\item
Describe a space~\(X\) with \(\pi(X)\cong\mathbb{Z}/3\mathbb{Z}\).
\item
If \(j:A\hookrightarrow X\) is the inclusion of a retract, then show that the induced map on fundamental groups \(j_\#\) is injective.
\item
Give the definition of a deformation retraction.
\item
Define Stone–Čech compactification.
\end{examproblems}
\end{document}
